\documentclass[11pt]{article}
\usepackage{amsmath,amssymb,amsthm,xcolor,geometry,hyperref,booktabs}
\geometry{margin=1in}
\newcommand{\chg}[1]{\textcolor{blue}{#1}}
\newcommand{\tr}{\operatorname{tr}}
\newcommand{\Res}{\operatorname{Res}}
\newcommand{\arctanh}{\operatorname{arctanh}}
\newcommand{\arccosh}{\operatorname{arccosh}}
\newcommand{\Pf}{\operatorname{Pf}}
\title{Challenge 64: Conformal correlators \\ \large Solution (independent derivation)}
\date{}
\begin{document}\maketitle

\section{Data and strategy}
Ising CFT ($c=\frac12$): $\epsilon$ with $h=\bar h=\frac12$, $\sigma$ with $h=\bar h=\frac1{16}$, normalised by $\langle\epsilon\epsilon\rangle=|z|^{-2}$, $\langle\sigma\sigma\rangle=|z|^{-1/4}$, and $C_{\sigma\sigma\epsilon}=\frac12$. Write $z_i=x_i$ ($i=1,2,3$), $a=x_4$, $b=x_5$, and $G\equiv\langle\epsilon(z_1)\epsilon(z_2)\epsilon(z_3)\sigma(a)\sigma(b)\rangle$.

I compute $G$ in two independent ways: (A) from the free Majorana fermion in the presence of two spin fields, and (B) by bosonising two copies of the Ising model. Both give the same function; they are then evaluated at the two configurations.

\section{Method A: Majorana fermion with two spin fields}
The energy operator is a fermion bilinear, $\epsilon=i\psi\bar\psi$. Fusing $\sigma\times\sigma=1+\epsilon$ and then with three $\epsilon$'s shows that only the channel $\sigma\sigma\to\epsilon$ contributes (the identity channel would require $\epsilon^3\to1$, which is impossible). Hence $G=|F|^2$ for a single (holomorphic $\times$ antiholomorphic) block, where $F$ is the chiral correlator of three fermions in the Ramond background created by the two spin fields:
\begin{align}
u_i&\equiv\frac{\langle\psi(z_i)\sigma(a)\sigma(b)\rangle}{\langle\sigma(a)\sigma(b)\rangle}=\frac{1}{\sqrt2}\frac{\sqrt{a-b}}{\sqrt{(z_i-a)(z_i-b)}},\\
K_{ij}&\equiv\frac{\langle\psi(z_i)\psi(z_j)\sigma(a)\sigma(b)\rangle}{\langle\sigma(a)\sigma(b)\rangle}
=\frac{1}{2(z_i-z_j)}\left[\sqrt{\frac{(z_i-a)(z_j-b)}{(z_i-b)(z_j-a)}}+\sqrt{\frac{(z_i-b)(z_j-a)}{(z_i-a)(z_j-b)}}\right]
=\frac{A_{ij}}{2(z_i-z_j)D_iD_j},
\end{align}
with $D_i=\sqrt{(z_i-a)(z_i-b)}$ and $A_{ij}=(z_i-a)(z_j-b)+(z_i-b)(z_j-a)$. Here $K_{ij}$ is the Ramond-sector propagator (it reduces to $\frac{\sqrt{z/w}+\sqrt{w/z}}{2(z-w)}$ for $a\to0$, $b\to\infty$, has the correct residue $1/(z_i-z_j)$, and the correct $(z_i-a)^{-1/2}$ branch points), and $u_i$ is fixed by conformal covariance with the normalisation $|u_i|^2\langle\sigma\sigma\rangle=\frac12|a-b|^{3/4}/|z_i-a||z_i-b|=\langle\epsilon\sigma\sigma\rangle$, i.e.\ by $C_{\sigma\sigma\epsilon}=\frac12$. Wick's theorem in the Ramond sector (a Pfaffian in which the zero mode supplies the one-point functions $u_i$) gives for an odd number of fermions
\begin{equation}
F=(a-b)^{-1/8}\big(u_1K_{23}-u_2K_{13}+u_3K_{12}\big)=\frac{(a-b)^{3/8}}{2\sqrt2\,D_1D_2D_3}\,\mathcal B,\qquad
\mathcal B\equiv\frac{A_{23}}{z_2-z_3}-\frac{A_{13}}{z_1-z_3}+\frac{A_{12}}{z_1-z_2}.
\end{equation}
Therefore
\begin{equation}
G=|F|^2=\frac{|a-b|^{3/4}\,|\mathcal B|^2}{8\prod_{i=1}^3|(z_i-a)(z_i-b)|}.
\label{eq:GA}
\end{equation}
Checks: (i) $z_2\to z_3$: $\mathcal B\to 2(z_3-a)(z_3-b)/(z_2-z_3)$, so $G\to|z_{23}|^{-2}\langle\epsilon(z_1)\sigma\sigma\rangle$, correct; (ii) $z_3\to a$: $\mathcal B\to A_{12}/z_{12}$ and $G\to\frac{1}{2|z_3-a|}\,|K_{12}|^2|a-b|^{-1/4}=\frac{C_{\sigma\epsilon\sigma}}{|z_3-a|}\langle\epsilon\epsilon\sigma\sigma\rangle$, correct; (iii) $F$ is antisymmetric in the $\psi$'s and $G$ is single valued.

\section{Method B: bosonisation of Ising$\,\times\,$Ising}
Two copies of the Ising model form a Dirac fermion, i.e.\ a free boson $\phi$ with $\langle\phi(z)\phi(w)\rangle=-\ln|z-w|^2$. Matching two-point functions,
\begin{equation}
\epsilon_1\epsilon_2=\pm\,\partial\phi\,\bar\partial\phi,\qquad \sigma_1\sigma_2=\sqrt2\cos\tfrac{\phi}{2},
\end{equation}
since $\langle\partial\phi\bar\partial\phi(z)\,\partial\phi\bar\partial\phi(0)\rangle=|z|^{-4}=\langle\epsilon\epsilon\rangle^2$ and $2\langle\cos\frac\phi2(z)\cos\frac\phi2(0)\rangle=|z|^{-1/2}=\langle\sigma\sigma\rangle^2$. Hence
\begin{equation}
G^2=\pm\big\langle\textstyle\prod_{i=1}^3\partial\phi\bar\partial\phi(z_i)\;2\cos\tfrac{\phi(a)}2\cos\tfrac{\phi(b)}2\big\rangle
=\pm\big\langle\textstyle\prod_i\partial\phi\bar\partial\phi(z_i)\,e^{i\phi(a)/2}e^{-i\phi(b)/2}\big\rangle ,
\end{equation}
only charge-neutral terms surviving (the two neutral terms are equal by $\phi\to-\phi$). By Wick's theorem, with $\langle e^{i\phi(a)/2}e^{-i\phi(b)/2}\rangle=|a-b|^{-1/2}$, the holomorphic and antiholomorphic sectors factorise and
\begin{equation}
G^2=|a-b|^{-1/2}\,|H|^2,\qquad H=f_1f_2f_3+\mathcal G_{12}f_3+\mathcal G_{13}f_2+\mathcal G_{23}f_1,
\end{equation}
where $f_i=-\frac i2\big(\frac{1}{z_i-a}-\frac{1}{z_i-b}\big)=-\frac{i(a-b)}{2D_i^2}$ is the contraction of $\partial\phi(z_i)$ with the vertex operators and $\mathcal G_{ij}=-(z_i-z_j)^{-2}$ is the $\partial\phi\partial\phi$ propagator (the overall sign is fixed by $G^2\ge0$). Explicitly
\begin{equation}
H=\frac{i(a-b)}{8\,D_1^2D_2^2D_3^2}\,Q,\qquad Q\equiv(a-b)^2+4\Big(\frac{D_1^2D_2^2}{z_{12}^2}+\frac{D_1^2D_3^2}{z_{13}^2}+\frac{D_2^2D_3^2}{z_{23}^2}\Big),
\end{equation}
so that
\begin{equation}
G=\frac{|a-b|^{3/4}\,|Q|}{8\prod_i|(z_i-a)(z_i-b)|}.
\label{eq:GB}
\end{equation}
Comparing (\ref{eq:GA}) and (\ref{eq:GB}) requires $|\mathcal B|^2=|Q|$. Indeed one verifies the polynomial identity $\mathcal B^2=Q$, so the two methods agree exactly. 

\section{Evaluation}
\paragraph{Configuration (1):} $z_1=1+i$, $z_2=2$, $z_3=3$, $a=4$, $b=5$. Then $A_{ij}=2z_iz_j-9(z_i+z_j)+40$ gives $A_{23}=7$, $A_{13}=10-3i$, $A_{12}=17-5i$; $z_{23}=-1$, $z_{13}=-2+i$, $z_{12}=-1+i$; so
\begin{equation}
\mathcal B=-7-\frac{10-3i}{-2+i}+\frac{17-5i}{-1+i}=-7+\Big(\frac{23}{5}+\frac45 i\Big)-(11+6i)=-\frac{67+26i}{5},\qquad |\mathcal B|^2=\frac{67^2+26^2}{25}=\frac{1033}{5}.
\end{equation}
The denominator: $|(z_1-4)(z_1-5)|=|11-7i|=\sqrt{170}$, $|(z_2-4)(z_2-5)|=6$, $|(z_3-4)(z_3-5)|=2$, and $|a-b|=1$. Thus
\begin{equation}
G_{(1)}=\frac{1033/5}{8\cdot12\sqrt{170}}=\frac{1033}{480\sqrt{170}}\approx0.1650572570 .
\end{equation}
\paragraph{Configuration (2):} $z_1=1$, $z_2=2$, $z_3=3$, $a=4$, $b=5$: $A_{23}=7$, $A_{13}=10$, $A_{12}=17$, $\mathcal B=-7+5-17=-19$, the denominator is $8\cdot12\cdot6\cdot2=1152$, hence
\begin{equation}
G_{(2)}=\frac{361}{1152}\approx0.3133680556 .
\end{equation}
Cross-check with method B for (2): $f_1=\frac{i}{24}$, $f_2=\frac{i}{12}$, $f_3=\frac{i}{4}$, $\mathcal G_{12}=\mathcal G_{23}=-1$, $\mathcal G_{13}=-\frac14$, $H=-i\big(\frac1{1152}+\frac14+\frac1{48}+\frac1{24}\big)=-\frac{361\,i}{1152}$, so $G=|H|=\frac{361}{1152}$; and $Q=1+4(72+6+12)=361=\mathcal B^2$. For (1) method B gives numerically $0.16505725697486\ldots$ as well.

\section{Comparison with the provided solution}
The provided solution uses method A with the same normalisations and obtains the same values. I found no errors. \chg{Added here: an independent derivation by bosonisation, the exact identity $\mathcal B^2=Q$ connecting the two, and an independent numerical confirmation of both configurations.}
\end{document}